\documentclass[10pt,a4paper]{article}
\usepackage[T1]{fontenc}
\usepackage[utf8]{inputenc}
\usepackage[spanish,es-tabla]{babel}
\usepackage{lmodern}
\usepackage[margin=2cm,headheight=32pt]{geometry}
\usepackage{graphicx,booktabs,longtable,tabularx,array,amsmath}
\usepackage{xcolor,tikz}
\usetikzlibrary{positioning,arrows.meta,fit,backgrounds,babel}
\usepackage{enumitem}
\usepackage[section]{placeins}
\usepackage{float}
\usepackage{needspace}
\usepackage{fancyhdr}
\usepackage{microtype}
\usepackage{cite}
\usepackage{pgfgantt}
\usepackage[font=footnotesize,labelfont=bf,labelsep=period]{caption}
\usepackage{parskip}
\PassOptionsToPackage{hyphens}{url}
\usepackage[hidelinks]{hyperref}
\hypersetup{pdftitle={Evaluación comparativa del filtrado de dominios para control parental},pdfauthor={Matías Vigneau Andrades},pdfsubject={Práctica I -- Universidad Diego Portales}}
\usepackage{listings}
\definecolor{udpblue}{HTML}{244A64}
\definecolor{softblue}{HTML}{EDF3F7}
\definecolor{softorange}{HTML}{FFF1DD}
\pagestyle{fancy}
\fancyhf{}
\fancyhead[L]{\includegraphics[width=1.2cm]{udp_logo.png}}
\fancyhead[R]{Facultad de Ingeniería y Ciencias}
\fancyfoot[R]{\thepage\hspace{0.02cm}}
\renewcommand{\headrulewidth}{0.9pt}
\renewcommand{\footrulewidth}{0.5pt}
\setlist{itemsep=3pt,topsep=5pt}
\setlength{\emergencystretch}{2em}
\renewcommand{\arraystretch}{1.15}
\newcolumntype{Y}{>{\raggedright\arraybackslash}X}
\lstset{basicstyle=\ttfamily\small,breaklines=true,frame=single,columns=fullflexible,keepspaces=true}

\begin{document}
\hypersetup{pageanchor=false}
\begin{titlepage}
\begin{center}
\vspace*{\baselineskip}
{\bfseries\fontsize{16}{19}\selectfont UNIVERSIDAD DIEGO PORTALES\\[0.5cm]}
{\fontsize{11}{13}\selectfont FACULTAD DE INGENIERÍA Y CIENCIAS\par}
\vspace*{0.5\baselineskip}
{\bfseries\fontsize{10}{12}\selectfont ESCUELA DE INFORMÁTICA Y TELECOMUNICACIONES\\[0.35cm]}
\vspace*{\baselineskip}
\includegraphics[scale=0.40]{udp_logo.png}\par
\vspace*{3\baselineskip}
\hrule height 0.5pt
\vspace{1mm}
\hrule height 1.5pt
\vspace*{\baselineskip}
{\bfseries\fontsize{14}{17}\selectfont Avances de Práctica I\\[0.5cm]}
{\bfseries\itshape\fontsize{12}{15}\selectfont Evaluación comparativa del filtrado de dominios para control parental\\en servicios de DNS protector, VPN y proxies\par}
\vspace*{\baselineskip}
\hrule height 0.5pt
\vspace{1mm}
\hrule height 1.5pt
\vspace*{4.5\baselineskip}
{\bfseries\fontsize{14}{17}\selectfont Profesor Guía: }{\fontsize{12}{15}\selectfont Nicolás Boettcher}\\[0.3cm]
{\bfseries\fontsize{14}{17}\selectfont Alumno: }{\fontsize{12}{15}\selectfont Matías Vigneau Andrades}\\[0.3cm]
\vspace{0.8cm}
{\fontsize{10}{12}\selectfont 30 de octubre de 2026\par}
\end{center}
\end{titlepage}
\tableofcontents
\thispagestyle{empty}
\clearpage
\setcounter{page}{1}
\hypersetup{pageanchor=true}

\section{Antecedentes y Motivación}
Este avance presenta la preparación del estudio y un piloto de consultas DNS a dos perfiles de Cloudflare. Las mediciones bajo VPN o proxy y la comparación completa de categorías forman parte del trabajo pendiente.

El Sistema de Nombres de Dominio permite consultar información asociada a nombres, por ejemplo las direcciones IPv4 e IPv6 mediante registros A y AAAA. La resolución recursiva obtiene respuestas de servidores autoritativos o reutiliza información almacenada temporalmente. Su funcionamiento se especifica en los RFC 1034 y 1035 \cite{rfc1034,rfc1035}. La finalidad básica de un resolutor es resolver nombres; la protección por categorías constituye una función adicional que debe identificarse en cada servicio.

En este trabajo se utiliza el término \emph{DNS protector} o PDNS para los servicios que aplican políticas destinadas a impedir la resolución ordinaria de ciertos dominios. El estudio de Liu y colaboradores describe proveedores con distintas formas de modificación de respuestas y distintas fuentes de información sobre amenazas \cite{liu2024}. Por ello, el nombre comercial de un servicio no basta para describir el tratamiento aplicado a cada consulta: también importan el perfil, la dirección consultada y la política activa.

Se llama \emph{perfil} a una configuración concreta de un servicio, con una dirección y una política de filtrado identificadas. Un ejemplo es Cloudflare, que documenta un perfil de seguridad y otro que agrega contenido adulto, con direcciones diferentes. Quad9 publica tanto una variante con protección contra amenazas como otra sin ese bloqueo \cite{cloudflare,quad9}. Para compararlas hay que conocer qué pretende bloquear cada perfil y qué información permite obtener su respuesta. Que un servicio de seguridad resuelva un sitio de una categoría ajena a su política no demuestra por sí solo una falla.

Las VPN y los proxies se consideran condiciones de conexión que pueden cambiar el recorrido de las solicitudes y el lugar de resolución. No ofrecen filtrado por el solo hecho de transportar tráfico. Quad9 advierte que un cliente VPN puede utilizar un resolutor distinto del configurado en el sistema \cite{quad9faq}. HTTP y SOCKS permiten distintas formas de intermediación \cite{rfc9110,rfc1928}. Por ello, su inclusión en la comparación requiere identificar una función de protección y comprobar qué tráfico utiliza esa conexión; una extensión del navegador no permite asumir el mismo recorrido para un programa de consola.

El foco es el filtrado de dominios para control parental: contrastar la política ofrecida con las respuestas que recibe el usuario, incluyendo el bloqueo indebido de contenido permitido. En el contexto chileno, Rojas, Boettcher, Palacios Játiva y Gutiérrez Gaitán proponen una métrica para evaluar las funciones de aplicaciones de control parental \cite{rojas2025}. Este trabajo aporta un antecedente de comprobación de las funciones ofrecidas por esas aplicaciones. La presente práctica aborda esa necesidad de evaluación desde las políticas y respuestas de filtrado de dominios. Por ello, utiliza un diseño de medición específico para los servicios DNS, sin trasladar directamente los datos o la métrica de las aplicaciones.

\textbf{¿Qué diferencias y qué estabilidad presentan las respuestas de perfiles de protección ante un corpus de contenido documentado para control parental, y qué evidencia permite distinguir señales de filtrado de fallas o respuestas negativas legítimas?}

Se distinguen contenido adulto, apuestas y armas de la detección de phishing o malware. Son criterios distintos: una página legítima puede pertenecer a una categoría que una política parental restringe. Los datos históricos de phishing se conservan para auditoría y no forman parte del corpus parental definitivo en este diseño. Se incluirán sitios educativos para estudiar bloqueos indebidos, después de revisar su contenido y comprobar que la política del servicio los permite. Que un sitio sea popular no sustituye esa revisión.

Para comparar los servicios también es necesario conocer la procedencia de las listas y el significado de sus categorías. Plonka, Palacio y Perouli estudian listas de amenazas frente a consultas observadas en una red de investigación y educación \cite{plonka2025}. Las diferencias de objetivos y procedencia que describen motivan documentar esas características en el corpus de esta práctica. Su estudio utiliza tráfico observado y listas de amenazas; la comparación parental requiere seleccionar y revisar dominios según su contenido.

La práctica reúne y compara respuestas de servicios de protección, prepara un corpus (conjunto de dominios de prueba) con fuentes y revisiones registradas y documenta cómo repetir las mediciones. El trabajo incluye revisar estudios previos, comprobar las categorías, configurar las conexiones, desarrollar el programa y analizar los resultados y sus límites.

Se compararán perfiles con políticas diferentes, cuya selección se justifica por sus funciones y no solo por su precio. Cloudflare permite comparar resolución ordinaria, seguridad y protección familiar dentro de un proveedor; Quad9 permite contrastar protección y resolución sin bloqueo; Google aporta una referencia de resolución adicional \cite{cloudflare,quad9,google}. Estas funciones justifican la selección inicial de seis perfiles. Es una decisión del diseño de esta práctica, revisable según conectividad y calibración; no una cantidad prescrita por los estudios citados.

Las etiquetas heredadas describen popularidad, contenido adulto, apuestas, armas y reportes de phishing. Se conserva el inventario original para auditar su procedencia, pero el alcance parental separa las amenazas de las categorías de contenido. La etiqueta heredada de armas requiere especial revisión: el archivo contiene indicios de sitios de odio y violencia que no acreditan pertenencia a armas. La revisión debe establecer la categoría aplicable a cada entrada antes de utilizarla en una comparación.

Se adopta como meta de trabajo reunir \textbf{2.000 dominios únicos revisados} para delimitar la construcción y comparación del corpus. Es una decisión de planificación, no un tamaño que los antecedentes establezcan como suficiente. La selección es intencional y su composición se justificará por las categorías y políticas evaluadas. Se informará cuántos dominios hay por categoría, de dónde provienen, cuáles aparecen en varias categorías y cuáles están revisados. Cada dominio añadido debe tener una función justificada en la comparación y conservar la evidencia de su selección. No se completará la cantidad con duplicados ni se dará por correcta una categoría sin revisarla.

La selección de fuentes también requiere distinguir su finalidad. Tranco ofrece un ranking de popularidad, útil para obtener candidatos reproducibles, pero no una clasificación de contenido permitido \cite{tranco}. UT1 publica listas por categorías de contenido y una lista de referencias de bibliotecas universitarias \cite{ut1}. Esas etiquetas de origen constituyen puntos de partida que deben contrastarse con evidencia de contenido y con la política del perfil. La cantidad de entradas de una lista no acredita por sí sola su calidad ni su pertinencia para el contexto evaluado.

Se utiliza Python porque la biblioteca dnspython permite enviar consultas y acceder a los mensajes, registros y campos DNS \cite{dnspython}. Esa capacidad permite automatizar las consultas y conservar la evidencia requerida. La herramienta registra y compara respuestas; no implementa un clasificador nuevo ni reconstruye el algoritmo interno de un proveedor.

Los trabajos de PDNS y OONI fundamentan distinguir respuestas modificadas, anomalías y fallas, e informar frecuencias y cobertura \cite{liu2024,ooni}. Las tablas de contingencia y el acuerdo observado describen las coincidencias entre perfiles; Kappa y Jaccard ofrecen descriptores adicionales con límites que se desarrollan en metodología \cite{cohen1960,kosub2016}. Su elección se apoya en esos antecedentes, pero el acuerdo entre servicios no equivale a acierto respecto del contenido.

Como primera comprobación del procedimiento se definió un piloto de 50 dominios para comprobar la selección, los controles, el registro y el análisis. La cantidad permite recorrer el procedimiento a pequeña escala; no acredita representatividad ni sustituye la meta de 2.000 revisados. Si las categorías solo cuentan con evidencia preliminar, el piloto se identificará como exploratorio y no se utilizará para estimar acierto del filtrado.

La comparación responderá tres preguntas: (i) qué protege cada perfil según su política declarada; (ii) qué respuestas entrega bajo condiciones reproducibles; y (iii) en qué casos sus resultados pueden compararse con los de otros perfiles o con mediciones de conectividad. Se adopta un estudio descriptivo y comparativo. El acuerdo entre perfiles y el acierto respecto de una política se analizan por separado; una discrepancia no demuestra por sí sola ineficacia.

\section{Objetivos}
\subsection{General}
Comparar las respuestas de servicios de DNS protector para control parental y de conexiones mediante VPN o proxy con una función de protección comprobada. Se utilizará un conjunto de dominios con categorías documentadas y consultas de control, y se registrarán las condiciones para repetir las pruebas y los límites de los resultados.

\subsection{Específicos}
\begin{enumerate}
\item Revisar y documentar los mecanismos de resolución y protección, las políticas declaradas de los perfiles y los métodos de medición, incluyendo OONI.
\item Construir y revisar un conjunto de dominios, conservando sus fuentes, categorías, criterios de inclusión, evidencias de revisión e historial de cambios.
\item Implementar consultas A y AAAA con controles antes y después de la observación, repeticiones, identificación de errores y almacenamiento de mensajes y configuración.
\item Documentar y verificar el recorrido de consultas directas y bajo VPN; tratar por separado la accesibilidad HTTP mediante proxies y las señales DNS.
\item Describir frecuencias, cobertura, estabilidad y acuerdo por perfil y categoría, indicando los casos excluidos y el total de casos utilizado para calcular cada proporción.
\item Entregar el código, las listas, los resultados y los pasos para repetir y revisar las pruebas.
\end{enumerate}

\section{Alcances y Limitaciones}
La configuración inicial incluye seis perfiles para comparar resolución sin bloqueo, protección contra amenazas y filtrado de contenido adulto. Son un punto de partida. Se debe ampliar y justificar el inventario según las categorías y funciones efectivamente disponibles. El archivo \texttt{config/profiles.json} conserva dirección, política, referencia y fuente. Un perfil de amenazas no se evalúa como si prometiera bloquear apuestas o armas. Para evaluar acierto parental, se requiere una política que declare las categorías correspondientes; la lista actual no acredita todavía esa cobertura completa. Las políticas se revisaron documentalmente el 27 de octubre de 2026 \cite{cloudflare,quad9,google}.

La cobertura se comprobará por categoría antes de ampliar las mediciones. Entre los perfiles activos, solo Cloudflare familiar declara tratar contenido adulto; apuestas y armas quedan pendientes de un perfil que declare esas categorías y cuya configuración pueda comprobarse. CleanBrowsing documenta categorías configurables de adulto, apuestas y armas en sus planes de pago \cite{cleanbrowsingapi,cleanbrowsingcategories}. Se incorpora como candidato documental, todavía sin cuenta configurada ni calibración. El estudio de esas categorías exigirá registrar el plan, activar el filtro correspondiente y verificar su recorrido y controles. Mientras falte ese paso, las consultas de apuestas y armas describen respuestas DNS, pero no evalúan acierto parental. El piloto no sustituye esta ampliación.

\begin{table}[H]
\centering\small
\begin{tabularx}{\textwidth}{p{3.5cm}p{2.2cm}Y}
\toprule
Perfil & Dirección & Papel en el diseño\\
\midrule
Cloudflare estándar & 1.1.1.1 & Referencia de resolución\\
Google Public DNS & 8.8.8.8 & Referencia adicional; no filtro parental\\
Cloudflare seguridad & 1.1.1.2 & Perfil de protección contra amenazas\\
Cloudflare familiar & 1.1.1.3 & Protección con contenido adulto\\
Quad9 seguro & 9.9.9.9 & Protección contra amenazas\\
Quad9 sin bloqueo & 9.9.9.10 & Referencia dentro del mismo proveedor\\
\bottomrule
\end{tabularx}
\caption{Perfiles configurados; su existencia en el archivo no acredita conectividad ni ejecución completa.}
\end{table}

Las VPN se incorporarán cuando se pueda identificar el resolutor, el cliente, plan, función de protección, ubicación del servidor y ruta efectiva. No se fija de antemano una comparación de tres planes gratuitos contra tres de pago: una prueba temporal o un plan distinto puede cambiar varias funciones simultáneamente. Cada condición debe describirse y justificarse antes de compararla. No se requiere contratar servicios para realizar las pruebas con los perfiles públicos iniciales.

En proxies se registrarán el tipo de intermediario, resolución local o remota, protocolo del destino, verificación TLS y huellas de bloqueo disponibles. La herramienta incluye una prueba HTTP opcional y limitada, sin ejecutar el contenido recibido ni seguir redirecciones automáticamente. La documentación de OONI sirve para comparar métodos de medición, y sus formatos permiten intercambiar datos. Para ejecutar sus pruebas habrá que registrar la configuración y dónde se guardarán o publicarán los resultados. En este avance no se han generado ni publicado mediciones de OONI.

El alcance anterior tiene límites: el corpus es un conjunto experimental de candidatos, no una muestra probabilística de Internet. Las categorías pueden superponerse y las fuentes pueden contener errores. La popularidad no acredita benignidad, ni una lista de phishing demuestra que sus entradas sean recientes o sigan activas. Los perfiles pueden perseguir objetivos distintos y evolucionar con el tiempo.

La herramienta actual envía consultas DNS mediante UDP o TCP al puerto 53. No implementa mediciones DNS sobre HTTPS (DoH) ni DNS sobre TLS (DoT); por tanto, no compara su comportamiento de filtrado con el del transporte convencional. Las diferencias entre protección DNSSEC y cifrado del transporte se explican en el estado del arte \cite{rfc4033,rfc8484,rfc7858}. Una prueba bajo VPN deberá verificar si estas consultas siguen el túnel, sin asumir que reproducen el DNS utilizado por el navegador.

El experimento describe una ubicación y ventanas temporales concretas. La configuración del equipo, el ISP, el enrutamiento y la caché pueden influir. No se controla la base interna del proveedor ni se demuestra el algoritmo que produjo una decisión. La fase DNS no garantiza acceso HTTP, seguridad del contenido, cumplimiento judicial ni protección integral del usuario. Estas preguntas requerirían evidencia adicional y un diseño específico.

\section{Resultados Esperados}
El resultado principal será una comparación documentada de perfiles de protección para control parental, con las condiciones necesarias para interpretar sus respuestas. Se entregarán los siguientes productos:
\begin{enumerate}
\item Un inventario de servicios y perfiles con fuentes, categorías declaradas, plan, configuración y criterios de inclusión. Las configuraciones bajo VPN y proxy se describirán según su recorrido efectivo y su función de protección.
\item Un corpus de 2.000 dominios únicos revisados, con evidencia de categoría, procedencia y decisiones de inclusión o exclusión. La entrega distinguirá ese corpus final de los candidatos y de los archivos históricos, e indicará cuántos faltan si no se alcanza la meta.
\item Una herramienta reproducible de medición, con configuración versionada, registro de mensajes y errores, controles de resolución y documentación de ejecución.
\item Tablas por perfil, categoría y condición de red con estados observados, cobertura y estabilidad; matrices de acuerdo calculadas sobre casos comparables. Cuando exista una decisión de referencia validada por política, se informarán también aciertos y bloqueos indebidos con sus denominadores.
\item Un informe que relacione los antecedentes, las decisiones del diseño y los resultados obtenidos, y un conjunto de listas y registros que puedan compararse con mediciones de OONI bajo condiciones documentadas.
\end{enumerate}
Los entregables no presuponen que un servicio será mejor que otro ni que una comparación entre planes gratuitos y de pago producirá diferencias. El resultado podrá incluir coincidencias, discrepancias o condiciones sin evidencia suficiente, siempre identificadas como tales.

\section{Estado del Arte}
\subsection{Resolución DNS y señales de protección}
Una consulta contiene, entre otros campos, el nombre, tipo y clase solicitados. La respuesta incluye un código de resultado, banderas y secciones de respuesta, autoridad e información adicional. Se diferencia el código de respuesta de las banderas: una decisión de permitir o bloquear no es una bandera del protocolo. Por ejemplo, RD solicita recursión, RA indica que el servidor la ofrece y TC indica truncamiento del mensaje \cite{rfc1035}. Se registran para describir cómo se obtuvo la respuesta; TC puede requerir repetir la consulta por TCP. No se utilizan como etiquetas de contenido.

Una respuesta NXDOMAIN expresa que el nombre consultado no existe; la ausencia de un registro del tipo solicitado tampoco equivale necesariamente a inexistencia del nombre. El RFC 2308 especifica las respuestas negativas y su almacenamiento en caché \cite{rfc2308}. En consecuencia, convertir cualquier NXDOMAIN o ausencia de respuesta A en un bloqueo produciría falsos indicios, especialmente en listas antiguas o volátiles. Un tiempo de espera agotado describe una consulta sin respuesta dentro del límite fijado; no revela la intención del proveedor.

El tiempo de vida (TTL) limita la reutilización de un registro, pero una observación aislada no demuestra que la respuesta provenga de caché ni permite reconstruir su edad original. Consultar una dirección explícita permite elegir el resolutor sin controlar por ello su caché. El TTL se conserva como contexto de la observación y no como certificado de ausencia de caché. Ni TTL ni latencia son indicadores principales de calidad de categorización; el objeto de evaluación es la respuesta respecto del contenido y de la política, no un ranking de rapidez.

DNSSEC proporciona autenticación del origen de los datos e integridad; no cifra el contenido de las consultas \cite{rfc4033}. DoH y DoT protegen el transporte mediante HTTPS y TLS, respectivamente \cite{rfc8484,rfc7858}. Son propiedades distintas, por lo que no se presentan como mecanismos equivalentes. Solicitar registros DNSSEC mediante el bit DO o recibir la bandera AD, que declara datos autenticados por el resolutor, no sustituye a una validación independiente realizada por el cliente \cite{rfc4033}. Esta distinción permite delimitar qué observa cada medición, independientemente del transporte empleado.

La modificación de una respuesta DNS puede consistir, por ejemplo, en sustituir la dirección habitual por una dirección destinada a impedir el acceso. Este es el significado de \emph{reescritura de la respuesta} en este informe. No se refiere a modificar el protocolo ni el nombre del servicio. Para relacionar una respuesta con una política deben distinguirse los siguientes niveles:
\begin{enumerate}
\item \textbf{Categoría de referencia del corpus:} etiqueta sustentada por una fuente y una revisión documentada.
\item \textbf{Política del perfil:} categorías o amenazas que el proveedor declara tratar.
\item \textbf{Respuesta observada:} datos del mensaje recibido y condiciones de la consulta.
\item \textbf{Interpretación:} señal de política, anomalía, resolución observada o caso indeterminado.
\end{enumerate}
Estos niveles se relacionan, pero ninguno sustituye automáticamente a los demás. Una decisión binaria no basta para reconstruir un algoritmo propietario.

Los códigos de error extendido DNS (EDE) agregan información sobre causas de respuesta, incluidas decisiones de política \cite{rfc8914}. Los códigos 15, 16 y 17 describen bloqueo, censura o filtrado. El código 18 puede describir una consulta de un cliente no autorizado y no basta para atribuir filtrado de contenido; los EDE también pueden acompañar respuestas con direcciones válidas. Su ausencia tampoco acredita que no exista filtrado. Las direcciones sumidero se reconocen con evidencia documentada y acotada al perfil: una especificación del proveedor o una calibración reproducible con controles conocidos. Se registra cuál de estas fuentes sustenta cada huella; una observación local no equivale a una garantía del proveedor. Por ejemplo, Cloudflare documenta la respuesta IPv4 0.0.0.0 para los dominios que bloquea \cite{cloudflare}. Las direcciones loopback, como 127.0.0.1, y las direcciones no especificadas tienen funciones distintas y deben registrarse por separado.

\subsection{Recorrido de una consulta: DNS, VPN y proxy}
Una VPN transporta tráfico mediante un túnel y puede asignar un resolutor o incorporar protección adicional. Ello depende del cliente y su configuración. La documentación de Quad9 explica que los clientes VPN pueden utilizar resolutores distintos de los configurados en el sistema \cite{quad9faq}. El experimento debe comprobar qué dirección recibe la consulta y qué ruta utiliza el equipo.

Un intermediario limitado al navegador puede afectar solamente a las solicitudes de ese navegador. Una consulta enviada por un programa de consola puede seguir otro recorrido. También debe diferenciarse un proxy HTTP de uno SOCKS. Si el cliente entrega el nombre del destino al proxy, este puede resolverlo; si entrega una dirección que ya resolvió, la consulta DNS precedió a esa solicitud \cite{rfc9110,rfc1928}. La figura \ref{fig:recorrido} distingue estas alternativas, que deben verificarse en cada cliente. Las observaciones de accesibilidad HTTP deben distinguirse de las respuestas DNS y de las decisiones del intermediario. Un estado HTTP 403 o una falla TLS no identifica por sí solo una política del intermediario.

\begin{figure}[H]
\centering
\begin{tikzpicture}[>=Stealth,font=\small,
 box/.style={draw=udpblue,rounded corners,fill=softblue,align=center,minimum height=1cm,text width=3.1cm,execute at begin node={\hyphenpenalty=10000\exhyphenpenalty=10000}},
 note/.style={font=\footnotesize,align=center,text width=3.4cm,anchor=north,execute at begin node={\hyphenpenalty=10000\exhyphenpenalty=10000}}]
\node[box] (a) at (0,0) {Programa Python\\consultas A y AAAA};
\node[box] (b) at (4.9,0) {Red del equipo\\ruta observada};
\node[box] (c) at (9.8,0) {Resolutor o PDNS\\perfil seleccionado};
\draw[->] (a)--(b);
\draw[->] (b)--(c);
\node[box] (d) at (0,-2.2) {Programa Python\\prueba bajo VPN};
\node[box,fill=softorange] (e) at (4.9,-2.2) {Interfaz del túnel\\ruta comprobada};
\node[box] (f) at (9.8,-2.2) {Resolutor asignado\\o explícito};
\draw[->] (d)--(e); \draw[->] (e)--(f);
\node[box] (g) at (0,-4.4) {Cliente HTTP / navegador\\DNS local (opcional)};
\node[box,fill=softorange] (h) at (4.9,-4.4) {Proxy HTTP / SOCKS\\DNS remoto si recibe nombre};
\node[box] (i) at (9.8,-4.4) {Destino web\\o respuesta del proxy};
\draw[->] (g)--node[midway,above,font=\scriptsize,fill=white,inner sep=1pt]{IP o nombre}(h); \draw[->] (h)--(i);
\node[note] at (0,-5.4) {La consola puede seguir\\una ruta distinta\\del navegador.};
\node[note] at (4.9,-5.4) {Comprobar dónde se resuelve\\el nombre y qué tráfico\\se transporta.};
\node[note] at (9.8,-5.4) {Distinguir la resolución DNS\\de la accesibilidad HTTP.};
\end{tikzpicture}
\caption{Recorridos que deben verificarse. Son esquemas experimentales, no una afirmación de que toda VPN o proxy incorpore filtrado.}
\label{fig:recorrido}
\end{figure}

\Needspace{8\baselineskip}
\subsection{Estudios previos y criterios de comparación}
Cohen formuló una medida de acuerdo para clasificaciones nominales que compara la coincidencia observada con la coincidencia esperada a partir de las proporciones de cada categoría en ambos evaluadores \cite{cohen1960}. La similitud de Jaccard compara la intersección y la unión de dos conjuntos \cite{kosub2016}; aquí permite describir la coincidencia de señales de bloqueo sin sumar como coincidencias los dominios que ambos perfiles resuelven. El acuerdo observado y Kappa responden a una pregunta de coincidencia entre secuencias de decisiones, no a si esas decisiones son correctas. Dos servicios pueden coincidir al resolver un dominio o al emitir una señal de bloqueo sin que ello identifique sus bases de datos, su proveedor ascendente o su algoritmo interno.

Para evaluar una política parental hace falta una referencia independiente: qué contenido tiene el dominio y qué decisión debería aplicar el perfil según su política declarada. Los contrastes de contenido permitido y restringido permiten distinguir coincidencia de acierto. Una lista de popularidad no aporta por sí sola esa clasificación de referencia comprobada. Por eso se conservarán los errores y los casos dudosos, y se indicará cuántos casos se usaron para calcular cada proporción. No todas las respuestas permiten decidir entre bloqueo y permiso.

Liu y colaboradores estudian el despliegue y las respuestas de PDNS mediante mediciones activas \cite{liu2024}. Su análisis de los formatos de respuesta orienta el registro de mensajes y la comprobación de señales de filtrado en esta práctica. Plonka y colaboradores estudian listas con diferentes objetivos mediante mediciones pasivas \cite{plonka2025}; su aporte es examinar la procedencia de las listas y las diferencias que dificultan compararlas. Ambos trabajos proporcionan antecedentes para el diseño, pero sus objetos de estudio y condiciones de medición difieren de la comparación parental propuesta.

OONI ofrece experimentos y documentación para interpretar mediciones de conectividad. Distingue bloqueos confirmados mediante huellas de las anomalías que requieren examinar causas alternativas \cite{ooni}. Su enfoque aporta una referencia para conservar controles y evitar interpretar toda falla como censura. Los resultados de \texttt{Web Connectivity} requieren conservar su contexto experimental; una observación DNS aislada no tiene necesariamente la misma unidad de análisis.

\begin{table}[H]
\centering\small
\begin{tabularx}{\textwidth}{p{2.4cm}YYY}
\toprule
Antecedente & Objeto & Utilidad para la práctica & Límite de su aplicación\\
\midrule
Liu et al. & PDNS y respuestas modificadas & Registro de respuestas & Estudia PDNS abiertos, no aplicaciones parentales\\
Plonka et al. & Listas y tráfico observado & Procedencia y objetivos de listas & Utiliza mediciones pasivas y listas de amenazas\\
OONI & Conectividad y anomalías & Controles y contraste independiente & Anomalía no equivale a bloqueo confirmado\\
Tranco & Ranking de popularidad & Candidatos de alta popularidad & Popularidad no garantiza benignidad\\
\bottomrule
\end{tabularx}
\caption{Aporte y límites de los antecedentes seleccionados.}
\end{table}

\section{Metodología de Trabajo}
\subsection{Selección y revisión de dominios}
El punto de partida son los cinco CSV heredados. Se conservan para auditar el trabajo previo. La construcción del corpus normaliza nombres mediante IDNA, rechaza IP, URL y etiquetas inválidas, agrupa dominios repetidos y conserva todas las etiquetas de origen. Cada fila del corpus registra nombre, categorías, archivos y líneas de procedencia, resúmenes SHA-256 para identificar los archivos de origen, estado de validación, evidencia, responsable y fecha de revisión.

Las etiquetas del inventario inicial son \emph{popular}, \emph{adultos}, \emph{apuestas}, \emph{armas} y \emph{phishing}. El corpus parental definitivo admite contenido adulto, apuestas, armas y contenido educativo revisado; los candidatos populares requieren su propia revisión. Phishing queda separado del alcance principal. No se afirma que las categorías sean excluyentes ni que todas las etiquetas recibidas sean correctas. El grupo histórico llamado «recientes» procede de una lista de phishing y no acredita antigüedad de registro. La categoría de control técnico se reserva para calibración y no integra las estadísticas de contenido.

\begin{table}[H]
\centering\small
\begin{tabularx}{\textwidth}{>{\raggedright\arraybackslash}p{2.9cm}YY}
\toprule
Fuente considerada & Utilidad & Comprobación necesaria\\
\midrule
Tranco & Nombres populares y referencia reproducible \cite{tranco} & ID de lista, fecha, política de selección; no equiparar a benignidad\\
UT1 & Categorías de contenido declaradas por la fuente \cite{ut1} & Copia fechada, documentación y revisión de etiquetas; fuente histórica incompleta\\
StevenBlack y selección manual & Candidatos relacionados con apuestas \cite{stevenblack} & Procedencia de cada subconjunto y contraste con otra evidencia\\
Phishing.Database & Fuente del inventario histórico de phishing \cite{phishingdatabase} & Estado, fecha y pertinencia de la señal; no inferir antigüedad\\
FINSIN & Posible contraste del inventario histórico \cite{finsin2022} & Disponibilidad, condiciones de acceso y actualidad\\
\bottomrule
\end{tabularx}
\caption{Fuentes consideradas para el corpus parental y la auditoría histórica; sus etiquetas requieren revisión.}
\end{table}

El conjunto candidato parental se guarda en \texttt{dataset/corpus-parental-candidatos.csv}: 400 adultos, 400 apuestas, 266 armas, 400 populares y 534 referencias de biblioteca, sin solapamientos entre estos grupos tras retirar la etiqueta phishing. La composición es una base de trabajo, no una cuota óptima ni una clasificación ya comprobada. Las entradas heredadas y las nuevas referencias comienzan en estado \texttt{pending}. Para pasar a \texttt{validated} se requiere justificar la etiqueta mediante evidencia identificable y revisión con responsable y fecha; las entradas descartadas quedan como \texttt{excluded}. El programa exige esos campos y conserva la decisión, pero no puede verificar por sí solo la calidad de la revisión humana. La categoría de cada dominio se revisa con evidencias independientes de las respuestas DNS y del acuerdo entre proveedores.

Para cada fuente se documentarán las definiciones de sus categorías, la fecha de la copia y el método publicado de asignación y actualización de etiquetas: revisión manual, procedimiento automático o combinación de ambos. Si la fuente no explica ese método, se registrará como información no disponible. Se comprobará además si dos listas reutilizan los mismos datos, para no tratar su coincidencia como evidencia independiente. Esta revisión de fuentes aún está pendiente.

La revisión de cada dominio conservará la URL o el documento consultado, fecha, una descripción breve de la evidencia y la razón de inclusión o exclusión. Se contrastará el contenido descrito con la definición de la categoría, sin deducirlo del nombre del dominio o de la respuesta DNS. Las discrepancias quedarán pendientes hasta resolverlas; si no hay evidencia suficiente, la entrada no pasará al corpus revisado. La actividad del sitio se registrará aparte de su categoría, sin exigir visitas automáticas a los sitios. Antes de medir se revisará la composición resultante y se reemplazarán los descartados conservando su historial.

La meta final es 2.000 dominios únicos que cumplan los criterios documentados. Si no se reúnen, se indicará cuántos faltan y el conjunto no se presentará como terminado. La automatización sirve para normalizar, auditar y reproducir la selección, no para reemplazar esa validación. Cuando se limite una ejecución, la herramienta selecciona sin reemplazo mediante una semilla registrada; ello hace reproducible la selección dentro del corpus, pero no la transforma en una muestra aleatoria de toda la web.

\subsection{Configuración y medición}
DNS admite TCP además de UDP; el RFC 7766 especifica los requisitos de su utilización \cite{rfc7766}. Antes de medir se comprobarán ambos transportes cuando la conectividad lo requiera. Cada ejecución registrará cuál utiliza, su tiempo de espera y sus controles, sin atribuir automáticamente una diferencia de respuesta al proveedor.

Antes del experimento se verifican casos técnicos conocidos: un dominio de ejemplo, los dominios de prueba publicados por Cloudflare y un nombre bajo el TLD reservado \texttt{.invalid} \cite{cloudflare,rfc2606,iana}. Se comprueba que el programa conserva respuestas y errores y que una respuesta negativa legítima no se etiqueta automáticamente como bloqueo. Se utilizaron \texttt{example.com} como control de resolución, \texttt{no-existe-pdns-lab.invalid} como respuesta negativa legítima, y \texttt{malware.testcategory.com} y \texttt{nudity.testcategory.com} como controles de las funciones publicadas por Cloudflare. Son distintos de los controles temporales antes y después de cada consulta; prueban el funcionamiento técnico y las señales de categoría, no la validación del corpus.

Cada ejecución recibe un identificador único, copia de la configuración y del corpus seleccionado, versiones, parámetros y huellas de los archivos de entrada. Para una VPN se exige una descripción local de la condición y se comprueba que la interfaz esté activa y que la ruta al resolutor declarado la utilice. Esa comprobación no demuestra por sí sola que la política protectora esté habilitada: su estado y evidencia de configuración se documentan aparte. Una extensión de navegador se prueba dentro de su ámbito; no se equipara a un túnel del sistema.

Una vez comprobadas esas condiciones, cada observación se identifica por \emph{dominio, perfil, tipo DNS, ronda y condición de red}. Se planifican tres rondas por ejecución para comprobar si las respuestas se mantienen. Ese número de repeticiones no garantiza por sí solo resultados estadísticamente representativos. En cada ronda se alterna el orden de perfiles mediante una semilla registrada. Los controles anterior y posterior consultan el mismo dominio y tipo al resolutor de referencia, antes y después de la consulta al perfil evaluado. En la configuración de seis perfiles, Google Public DNS, Cloudflare seguridad y Cloudflare familiar usan 1.1.1.1; Quad9 seguro usa 9.9.9.10. Cloudflare estándar y Quad9 sin bloqueo se consultan sin esos controles. En el piloto de dos perfiles solo Cloudflare familiar utiliza controles, con 1.1.1.1 como referencia. Realizar estas consultas en momentos cercanos permite detectar cambios o fallas que dificultarían la comparación.

Se consultan A y AAAA por separado, solicitando información DNSSEC disponible. Se conserva el mensaje enviado y recibido, hora UTC, direcciones y CNAME, códigos de respuesta y EDE, TTL, banderas, transporte y duración. UDP utiliza respaldo TCP ante truncamiento. Los errores de red pueden reintentarse una vez por defecto y cada intento queda registrado. Una respuesta negativa no se reemplaza por reintentos hasta obtener el resultado deseado. El intervalo configurable entre consultas comienza en 0,2 segundos; es una decisión operativa revisable, no una garantía contra límites de consulta del proveedor.

Para $N$ dominios, $S$ perfiles, $T$ tipos y $R$ rondas, el número de observaciones del perfil es $NSTR$. Si $S_c$ perfiles requieren dos controles por observación, los grupos de consulta son $NTR(S+2S_c)$, antes de reintentos. Con seis perfiles, dos tipos, tres rondas y cuatro perfiles con controles, corresponden $36N$ observaciones y $84N$ grupos de consulta. La opción de simulación informa estos valores antes de transmitir consultas; la carga se ajustará tras la calibración.

\begin{figure}[H]
\centering
\begin{tikzpicture}[>=Stealth,font=\small,
 box/.style={draw=udpblue,rounded corners,fill=softblue,align=center,text width=3.5cm,minimum height=1cm}]
\node[box] (a) at (0,0) {Copias de las fuentes\\procedencia y SHA-256};
\node[box] (b) at (5.1,0) {Lista de dominios\\validación de etiquetas};
\node[box] (c) at (10.2,0) {Plan experimental\\perfiles y contexto};
\node[box] (d) at (0,-2.2) {Control anterior\\consulta de referencia};
\node[box,fill=softorange] (e) at (5.1,-2.2) {Perfil evaluado\\mensaje completo};
\node[box] (f) at (10.2,-2.2) {Control posterior\\consulta de referencia};
\node[box] (g) at (5.1,-4.4) {SQLite y análisis\\cobertura\\e incertidumbre};
\draw[->] (a)--(b);\draw[->] (b)--(c);
\draw[->] (c.south)--++(0,-0.35)-| (d.north);
\draw[->] (d)--(e);\draw[->] (e)--(f);
\draw[->] (f.south)--++(0,-0.4)-| (g.north);
\end{tikzpicture}
\caption{Etapas de preparación, medición y análisis. Los perfiles de referencia sin contraste utilizan solo la consulta al propio perfil.}
\label{fig:flujo}
\end{figure}

\subsection{Interpretación y análisis de resultados}
Las respuestas se clasifican después de guardar la evidencia. Se utiliza el término \emph{señalizado} para precisar que se observa una señal declarada o documentada, no una prueba absoluta sobre la intención ni una validación del contenido.
\begin{table}[H]
\centering\small
\begin{tabularx}{\textwidth}{p{3cm}Y}
\toprule
Estado & Criterio operativo\\
\midrule
Bloqueo señalizado & EDE de bloqueo (15--17) o sumidero con evidencia del proveedor o de calibración reproducible, identificada para ese perfil y configuración.\\
Resolución observada & NOERROR con direcciones de respuesta que no coinciden con una huella reconocida ni con direcciones loopback o no especificadas. No garantiza accesibilidad web.\\
Anomalía & Respuesta negativa con controles anterior y posterior que resuelven, o dirección especial sin huella que permita atribuirla. Requiere investigación adicional.\\
Error de red & Falla de transporte o tiempo agotado; se conservan sus intentos.\\
Indeterminado & Ausencia de evidencia suficiente, por ejemplo respuestas negativas también presentes en los controles.\\
\bottomrule
\end{tabularx}
\caption{Estados implementados; las anomalías y errores no se convierten en decisiones binarias de bloqueo.}
\end{table}

Las huellas deben comprobarse durante la calibración. Los mensajes DNS convencionales pueden verse afectados por intermediarios y los EDE no constituyen por sí mismos autenticación criptográfica. También pueden existir bloqueos no detectados por estas reglas. Estas reglas permiten revisar cómo se interpretó cada respuesta. La falta de una señal reconocida no descarta que exista filtrado.

A partir de estos estados se informa, por perfil, categoría y tipo DNS, el total de observaciones y la cantidad en cada estado. La proporción de señales sobre el total es distinta de la proporción entre observaciones interpretables. Ambas se acompañan de la cobertura: la proporción de observaciones que pueden interpretarse como señal de bloqueo o resolución. Los errores no se rellenan como permisos ni bloqueos. Los dominios con varias etiquetas pueden contribuir a varios desgloses, pero solo una vez al resumen global. No se suman desgloses solapados para reconstruir el total.

Para comparar perfiles se exige estabilidad del estado interpretable en las rondas previstas para cada dominio y tipo. Se informa cuántos dominios quedan fuera por cambio, error o cobertura incompleta. Cada pareja se compara con sus propios casos disponibles; la falta de datos de un tercer perfil no elimina esos casos. Se presenta la tabla de contingencia, el acuerdo observado y la similitud de Jaccard, $J=|B_1\cap B_2|/|B_1\cup B_2|$, donde $B_1$ y $B_2$ son los conjuntos señalados por cada perfil dentro de los mismos casos comparables \cite{kosub2016}. Si ninguno señala bloqueos, este informe registra la medida como no definida por falta de casos positivos; es una decisión de presentación del estudio, en lugar de adoptar la convención de similitud igual a uno para dos conjuntos vacíos.

El Kappa de Cohen se utiliza como medida adicional de acuerdo entre dos series de resultados codificados como 0 o 1 \cite{cohen1960}:
\[
\kappa=\frac{p_o-p_e}{1-p_e},\qquad
p_e=p_Ap_B+(1-p_A)(1-p_B).
\]
Aquí $p_o$ es la fracción de coincidencias y $p_A,p_B$ las proporciones de señales de bloqueo de cada perfil en los mismos casos. Se publican esas proporciones y el número de casos comparados. Cuando $p_e=1$, el valor se informa como no definido. Si un perfil nunca señala bloqueos y el otro sí, $p_o=p_e$ y Kappa es cero por construcción. Del mismo modo, Jaccard es cero si un conjunto de bloqueos está vacío y el otro no. En esos casos se conservan los valores por reproducibilidad, se advierte su falta de capacidad para distinguir comportamientos y se interpreta la tabla de contingencia. No se usa una escala universal para aceptar o rechazar la eficacia de un servicio.

La evaluación de acierto respecto de una política requerirá categorías validadas y criterios de decisión esperada por perfil, definidos antes del análisis. Se distinguirá bloqueo de contenido restringido, paso de contenido restringido, bloqueo indebido de contenido permitido y resolución de contenido permitido. Los errores e indeterminados se informarán aparte. La comparación incluirá candidatos educativos revisados cuando la política los permita, sin declarar benigno todo dominio popular. Solo en ese subconjunto corresponderá estimar medidas de acierto o errores de política, con sus denominadores. La versión actual produce descripciones de evidencia y acuerdo, no una clasificación global de proveedores de mejor a peor ni una estimación de falsos positivos a partir de popularidad.

Para el contraste con OONI se exportan a una lista de URL únicamente dominios con revisión documentada. La ejecución de OONI, si se realiza, debe conservar versión, hora, red, perfil y configuración, y tener en cuenta que sus resultados pueden publicarse según la configuración utilizada \cite{ooni}. Esta herramienta no ejecuta ni publica pruebas OONI automáticamente. El importador local conserva los valores originales de accesibilidad y bloqueo para revisarlos junto con el contexto de la medición.

Las mediciones OONI se contrastarán solo cuando coincidan dominio, período y condición de red de manera suficientemente documentada. Una diferencia puede originarse en capa, hora o punto de observación. Se registrará y analizará esa diferencia; no se impondrá la conclusión de una herramienta sobre la otra. Git mantiene código, fuentes del informe, configuración pública y auditorías; las mediciones detalladas se guardan por ejecución y se comparten mediante extractos revisados cuando corresponda.

\subsection{Plan de trabajo y criterios de cierre}
El trabajo se organiza en cinco etapas, descritas en la tabla \ref{tab:plan}. Para seleccionar los dominios definitivos hay que revisar sus categorías; antes de medir, hay que comprobar los perfiles y las rutas de las consultas. El plan organiza el trabajo completo, sin asignar fechas ni contabilizar horas trabajadas. El piloto descrito en desarrollo inicia la medición y el análisis a pequeña escala; la evaluación completa y la validación final de categorías siguen pendientes. Las fechas de ejecución y cierre quedan pendientes de confirmar; este esquema muestra dependencias, no una carta Gantt calendarizada.

\begin{table}[H]
\centering\small
\begin{tabularx}{\textwidth}{p{3.1cm}Y}
\toprule
Etapa & Criterio de cierre\\
\midrule
Fuentes y corpus & Revisar las categorías, incluir sitios de contenido permitido y documentar los 2.000 dominios únicos o cuántos faltan.\\
Configuración y calibración & Verificar conectividad, huellas de política y rutas; completar y justificar la selección de perfiles.\\
Medición & Completar las rondas previstas bajo condiciones registradas, conservando controles e incidencias.\\
Análisis y contraste & Publicar cobertura, estabilidad y acuerdo; contrastar políticas y OONI donde exista evidencia comparable.\\
Informe y revisión & Integrar fuentes, métodos y avances durante todas las etapas; redactar conclusiones al finalizar el análisis.\\
\bottomrule
\end{tabularx}
\caption{Etapas del trabajo y requisitos para completarlas.}
\label{tab:plan}
\end{table}

\begin{figure}[H]
\centering
\begin{ganttchart}[hgrid,vgrid,x unit=1.35cm,y unit chart=0.55cm,
 title label font=\small,bar label font=\small,bar/.style={fill=gray!55}]{1}{5}
\gantttitle{Orden de las tareas}{5}\\
\gantttitlelist{1,...,5}{1}\\
\ganttbar{Fuentes y corpus}{1}{1}\\
\ganttbar{Calibración de perfiles}{2}{2}\\
\ganttbar{Mediciones comparables}{3}{3}\\
\ganttbar{Análisis y contraste}{4}{4}\\
\ganttbar{Redacción y revisión}{1}{5}
\end{ganttchart}
\caption{Orden previsto de las tareas. El informe se redactará durante todo el proceso. Las columnas 1--4 corresponden a las primeras cuatro etapas de la tabla; la quinta tarea, redacción, acompaña el proceso completo. Las columnas no representan semanas.}
\end{figure}

\section{Desarrollo}
\subsection{Implementación y prueba piloto}
El repositorio es el espacio de trabajo del proyecto y un entorno virtual \texttt{.venv} contiene sus dependencias Python. La interfaz \texttt{pdns-lab} permite auditar el corpus, registrar revisiones, simular el plan, medir y analizar. El informe LaTeX se mantiene junto al código y se actualiza en el proyecto de Overleaf.

El módulo de planificación separa el inventario histórico, los candidatos en alcance y la selección final; exige la cantidad y validación declaradas en \texttt{config/estudio.json}. La estructura del proyecto separa \texttt{src/pdns\_lab/}, \texttt{config/}, \texttt{dataset/}, \texttt{tests/} y \texttt{docs/informe/}. Las mediciones nuevas se almacenan en \texttt{results/}, identificando cada ejecución por separado. La base histórica \texttt{resultados.db} se conserva como antecedente exploratorio; sus etiquetas anteriores no se copian como decisiones de bloqueo comprobadas. El código anterior se conserva en \texttt{legacy/} para consultar cómo se obtuvieron esos registros.

En una comprobación previa se consultaron cuatro dominios técnicos contra Cloudflare estándar y familiar, con consultas A, una ronda, un segundo de espera y sin reintentos. Se registraron ocho observaciones del perfil. Las cuatro consultas al perfil familiar agotaron el tiempo de espera desde este equipo; en el perfil estándar se observaron tres resoluciones y una respuesta negativa indeterminada del nombre reservado. Estos resultados verifican el registro y la separación de errores, pero son insuficientes para evaluar la protección del perfil familiar o atribuir la causa de la falta de respuesta.

Las pruebas automatizadas cubren respuestas negativas legítimas, controles incompletos, sumideros documentados, EDE, múltiples direcciones, reintentos, categorías solapadas, valores degenerados de Kappa, datos faltantes por par, rondas inestables, importación OONI y estados HTTP ambiguos. Sus respuestas de red son simuladas, por lo que verifican el comportamiento del programa y no la calidad de un proveedor. La evidencia de ejecución y los comandos de reproducción se conservan en el repositorio.


El 30 de octubre de 2026 se preparó un piloto de 50 dominios: 39 portales educativos o de investigación, cuatro plataformas de contenido adulto, cuatro sitios de apuestas y tres fabricantes de armas. La selección fue intencional, según la evidencia documental disponible, con predominio educativo. Se revisaron páginas institucionales, funciones de los sitios y catálogos oficiales; para adultos se utilizó una comunicación de la Comisión Europea que identifica las cuatro plataformas \cite{ec2025adult}. Se registraron por dominio la fuente, fecha, categoría propuesta y límites. Esa evidencia es preliminar: las filas siguen pendientes de validación final y la ejecución queda marcada como exploratoria. Los 2.000 candidatos principales conservan su estado anterior.

Se consultaron primero cuatro controles técnicos en los seis perfiles por UDP, para A y AAAA: 48 observaciones, de las cuales 28 terminaron en error de red. Una comprobación posterior por TCP recibió las 48 respuestas. Los tiempos de espera fueron 1,5 segundos en UDP y 2,5 en TCP, por lo que la diferencia no permite aislar una causa. El perfil familiar devolvió 0.0.0.0 y :: en los controles de categoría, mientras el estándar entregó direcciones de resolución. La dirección \texttt{::} se incorporó como huella empírica exclusiva de este piloto. Su fuente es el registro reproducible \texttt{docs/evidencia/calibracion-piloto-2026-10-30.json}: contiene dominio de control, perfil, tipo, direcciones y transporte. Esa procedencia se conserva en \texttt{fingerprint\_source}; no se presenta como una especificación publicada por Cloudflare ni se extiende a otros proveedores.

El piloto comparó Cloudflare estándar y familiar por TCP/53, con A y AAAA y tres rondas. Se registraron semilla 20261030, intervalo de 0,2 segundos, espera de 2,5 segundos y cero reintentos. La ruta local utiliza la interfaz inalámbrica; no se activó una VPN ni un proxy. Con los controles del perfil familiar se planificaron 600 observaciones y 1.200 consultas. La configuración completa se conserva en el repositorio.

El piloto completó las 600 observaciones previstas. La tabla \ref{tab:piloto} resume los estados después de corregir una regla detectada durante la revisión. El clasificador anterior interpretaba EDE 18 como bloqueo incluso en respuestas con direcciones válidas. Se retiró ese código de la regla automática conforme al RFC 8914 \cite{rfc8914}. La nueva interpretación cambió 36 estados; se generó en una base distinta, conservando los mensajes y la clasificación originales, el identificador de procedencia y las huellas SHA-256. Las 28 pruebas automatizadas, incluidas las regresiones de este caso y la conservación de evidencia, finalizaron correctamente.

\begin{table}[H]
\centering\small
\begin{tabular}{llrrrrr}
\toprule
Perfil & Tipo & Total & Resolución & Señal & Indeterminado & Error\\
\midrule
Estándar & A & 150 & 150 & 0 & 0 & 0\\
Estándar & AAAA & 150 & 30 & 0 & 120 & 0\\
Familiar & A & 150 & 138 & 12 & 0 & 0\\
Familiar & AAAA & 150 & 27 & 12 & 111 & 0\\
\bottomrule
\end{tabular}
\caption{Piloto exploratorio por TCP: observaciones en tres rondas, no dominios únicos.}
\label{tab:piloto}
\end{table}

En A, el perfil estándar resolvió los 50 dominios en las tres rondas. El familiar señaló los cuatro dominios con categoría preliminar de adulto y resolvió los otros 46. No se observaron señales sobre los 39 portales educativos en A; ello no estima una tasa general de bloqueos indebidos. La resolución de apuestas y armas no se interpreta como una falla parental: la política consultada del perfil familiar no promete bloquear esas categorías por sí solas. No hubo errores de red en el piloto TCP.

La comparación entre perfiles incluye 50 dominios para A: 46 coincidencias de resolución y cuatro diferencias, con acuerdo observado de 0,92. En AAAA solo diez dominios conservaron estados interpretables en ambos perfiles durante las tres rondas; hubo nueve coincidencias y una diferencia, con acuerdo 0,90. Los otros 40 se excluyen de esa comparación, principalmente por respuestas sin direcciones. No se convierten en permisos o bloqueos. El estándar no señala bloqueos en ninguno de los casos comparables. Por ello, Kappa y Jaccard son cero por construcción y no aportan una comparación discriminante en este piloto. La tabla \ref{tab:contingencia} muestra directamente los casos: las diferencias observadas son cuatro dominios en A y uno en AAAA, sin estimar su acierto respecto de una política validada.

\begin{table}[H]
\centering\small
\begin{tabular}{llrr}
\toprule
Tipo & Estándar / Familiar & Resolución & Señal\\
\midrule
A & Resolución & 46 & 4\\
A & Señal & 0 & 0\\
AAAA & Resolución & 9 & 1\\
AAAA & Señal & 0 & 0\\
\bottomrule
\end{tabular}
\caption{Contingencia por dominio estable. Filas: estado del estándar; columnas: estado del familiar. En AAAA se excluyen 40 dominios sin estados comparables.}
\label{tab:contingencia}
\end{table}

El piloto permitió completar consultas, controles y análisis, y corregir la interpretación a partir de los mensajes. La selección intencional, las categorías preliminares, una sola conexión y la cobertura reducida en AAAA impiden generalizar sus resultados. Faltan la validación final y las mediciones del diseño completo. El repositorio conserva la selección, configuración y extractos de resultados; los mensajes detallados quedan guardados localmente por ejecución.


\subsection{Auditoría y registros históricos}
La tabla \ref{tab:auditoria} muestra el resultado reproducible de la validación sintáctica. Los nombres de archivo conservan «top400» por procedencia, pero esa denominación no acredita la cantidad ni la calidad de su contenido.
\begin{table}[H]
\centering
\begin{tabular}{lrrr}
\toprule
Lista heredada & Filas & Rechazadas & Membresías válidas\\
\midrule
Adultos & 400 & 0 & 400\\
Armas & 360 & 94 & 266\\
Apuestas & 400 & 0 & 400\\
Phishing (antes «recientes») & 400 & 0 & 400\\
Popularidad & 400 & 0 & 400\\
\midrule
Total & 1.960 & 94 & 1.866\\
\bottomrule
\end{tabular}
\caption{Auditoría sintáctica de los CSV. Las 1.866 membresías corresponden a 1.852 dominios únicos por 14 solapamientos.}
\label{tab:auditoria}
\end{table}

Las 94 entradas rechazadas corresponden a 93 direcciones IPv4 y una etiqueta inválida en el archivo de armas. Aún falta revisar el contenido y la categoría de todas las entradas admitidas. Los 1.852 nombres describen el inventario histórico; no son los 2.000 dominios revisados de la meta actual ni el corpus parental definitivo. Los 14 solapamientos corresponden exclusivamente a dominios con etiquetas heredadas de adulto y phishing. Al retirar la etiqueta phishing se conservan como candidatos de adulto, por lo que los grupos del conjunto parental quedan sin solapamientos; su categoría sigue pendiente de revisión. El informe de auditoría conserva cada rechazo con archivo, línea y motivo. La auditoría permite conservar la diferencia entre cantidad de filas, dominios únicos y categorías revisadas.

La base histórica contiene 13.720 registros: 1.960 para cada uno de los cuatro resolutores públicos identificados como \texttt{Cisco\_Umbrella}, \texttt{Cloudflare\_Families}, \texttt{Google\_DNS} y \texttt{Quad9\_Security}, y 1.960 para cada una de las condiciones \texttt{ProtonVPN\_Free}, \texttt{Windscribe\_Free} y \texttt{HideMe\_Free}. Los recuentos se comprobaron mediante consultas de solo lectura. Son filas de ejecución sobre las listas originales, no 13.720 dominios únicos ni observaciones del corpus parental revisado.

Cisco\_Umbrella pertenece a la ejecución histórica y no a la configuración actual de seis perfiles; su nombre se conserva como etiqueta de procedencia. Esa etiqueta no acredita qué política estaba activa. FINSIN se revisó como posible fuente chilena de contraste de esos registros de amenazas \cite{finsin2022}, sin confirmar disponibilidad actual de sus datos.

Los archivos recuperados de Windows conservan condición, dominio, categoría, estado, dirección y hora, pero no acreditan la ruta, política activa ni controles del diseño actual. Se mantienen como antecedentes exploratorios. El piloto en Linux se guardó como una ejecución nueva, con configuración y controles propios. Las futuras condiciones VPN deberán documentar además cliente, plan, servidor e interfaz. La coincidencia entre registros históricos no demuestra que los servicios compartan infraestructura o algoritmo.

\section{Conclusiones}
% Se completará cuando termine la práctica y existan resultados finales.

\newpage
\begingroup
\footnotesize
\setlength{\parskip}{0pt}
\phantomsection
\addcontentsline{toc}{section}{Referencias}
\bibliographystyle{IEEEtran}
\bibliography{referencias}
\endgroup
\end{document}
